% Options for packages loaded elsewhere
% Options for packages loaded elsewhere
\PassOptionsToPackage{unicode}{hyperref}
\PassOptionsToPackage{hyphens}{url}
\PassOptionsToPackage{dvipsnames,svgnames,x11names}{xcolor}
%
\documentclass[
  letterpaper,
  DIV=11,
  numbers=noendperiod]{scrartcl}
\usepackage{xcolor}
\usepackage{amsmath,amssymb}
\setcounter{secnumdepth}{-\maxdimen} % remove section numbering
\usepackage{iftex}
\ifPDFTeX
  \usepackage[T1]{fontenc}
  \usepackage[utf8]{inputenc}
  \usepackage{textcomp} % provide euro and other symbols
\else % if luatex or xetex
  \usepackage{unicode-math} % this also loads fontspec
  \defaultfontfeatures{Scale=MatchLowercase}
  \defaultfontfeatures[\rmfamily]{Ligatures=TeX,Scale=1}
\fi
\usepackage{lmodern}
\ifPDFTeX\else
  % xetex/luatex font selection
\fi
% Use upquote if available, for straight quotes in verbatim environments
\IfFileExists{upquote.sty}{\usepackage{upquote}}{}
\IfFileExists{microtype.sty}{% use microtype if available
  \usepackage[]{microtype}
  \UseMicrotypeSet[protrusion]{basicmath} % disable protrusion for tt fonts
}{}
\makeatletter
\@ifundefined{KOMAClassName}{% if non-KOMA class
  \IfFileExists{parskip.sty}{%
    \usepackage{parskip}
  }{% else
    \setlength{\parindent}{0pt}
    \setlength{\parskip}{6pt plus 2pt minus 1pt}}
}{% if KOMA class
  \KOMAoptions{parskip=half}}
\makeatother
% Make \paragraph and \subparagraph free-standing
\makeatletter
\ifx\paragraph\undefined\else
  \let\oldparagraph\paragraph
  \renewcommand{\paragraph}{
    \@ifstar
      \xxxParagraphStar
      \xxxParagraphNoStar
  }
  \newcommand{\xxxParagraphStar}[1]{\oldparagraph*{#1}\mbox{}}
  \newcommand{\xxxParagraphNoStar}[1]{\oldparagraph{#1}\mbox{}}
\fi
\ifx\subparagraph\undefined\else
  \let\oldsubparagraph\subparagraph
  \renewcommand{\subparagraph}{
    \@ifstar
      \xxxSubParagraphStar
      \xxxSubParagraphNoStar
  }
  \newcommand{\xxxSubParagraphStar}[1]{\oldsubparagraph*{#1}\mbox{}}
  \newcommand{\xxxSubParagraphNoStar}[1]{\oldsubparagraph{#1}\mbox{}}
\fi
\makeatother


\usepackage{longtable,booktabs,array}
\usepackage{calc} % for calculating minipage widths
% Correct order of tables after \paragraph or \subparagraph
\usepackage{etoolbox}
\makeatletter
\patchcmd\longtable{\par}{\if@noskipsec\mbox{}\fi\par}{}{}
\makeatother
% Allow footnotes in longtable head/foot
\IfFileExists{footnotehyper.sty}{\usepackage{footnotehyper}}{\usepackage{footnote}}
\makesavenoteenv{longtable}
\usepackage{graphicx}
\makeatletter
\newsavebox\pandoc@box
\newcommand*\pandocbounded[1]{% scales image to fit in text height/width
  \sbox\pandoc@box{#1}%
  \Gscale@div\@tempa{\textheight}{\dimexpr\ht\pandoc@box+\dp\pandoc@box\relax}%
  \Gscale@div\@tempb{\linewidth}{\wd\pandoc@box}%
  \ifdim\@tempb\p@<\@tempa\p@\let\@tempa\@tempb\fi% select the smaller of both
  \ifdim\@tempa\p@<\p@\scalebox{\@tempa}{\usebox\pandoc@box}%
  \else\usebox{\pandoc@box}%
  \fi%
}
% Set default figure placement to htbp
\def\fps@figure{htbp}
\makeatother


% definitions for citeproc citations
\NewDocumentCommand\citeproctext{}{}
\NewDocumentCommand\citeproc{mm}{%
  \begingroup\def\citeproctext{#2}\cite{#1}\endgroup}
\makeatletter
 % allow citations to break across lines
 \let\@cite@ofmt\@firstofone
 % avoid brackets around text for \cite:
 \def\@biblabel#1{}
 \def\@cite#1#2{{#1\if@tempswa , #2\fi}}
\makeatother
\newlength{\cslhangindent}
\setlength{\cslhangindent}{1.5em}
\newlength{\csllabelwidth}
\setlength{\csllabelwidth}{3em}
\newenvironment{CSLReferences}[2] % #1 hanging-indent, #2 entry-spacing
 {\begin{list}{}{%
  \setlength{\itemindent}{0pt}
  \setlength{\leftmargin}{0pt}
  \setlength{\parsep}{0pt}
  % turn on hanging indent if param 1 is 1
  \ifodd #1
   \setlength{\leftmargin}{\cslhangindent}
   \setlength{\itemindent}{-1\cslhangindent}
  \fi
  % set entry spacing
  \setlength{\itemsep}{#2\baselineskip}}}
 {\end{list}}
\usepackage{calc}
\newcommand{\CSLBlock}[1]{\hfill\break\parbox[t]{\linewidth}{\strut\ignorespaces#1\strut}}
\newcommand{\CSLLeftMargin}[1]{\parbox[t]{\csllabelwidth}{\strut#1\strut}}
\newcommand{\CSLRightInline}[1]{\parbox[t]{\linewidth - \csllabelwidth}{\strut#1\strut}}
\newcommand{\CSLIndent}[1]{\hspace{\cslhangindent}#1}



\setlength{\emergencystretch}{3em} % prevent overfull lines

\providecommand{\tightlist}{%
  \setlength{\itemsep}{0pt}\setlength{\parskip}{0pt}}



 


% ============================================================
% loani.tex  —  include-in-header snippet, LuaLaTeX
% IPA handled by ipa.lua Lua filter at AST level.
% No ucharclasses, no babel, no runtime font switching.
% ============================================================

\usepackage{iftex}
\RequireLuaTeX

\usepackage{fontspec}

\setmainfont{Georgia}[
  Ligatures = TeX
]
% \usepackage[hebrew]{babel}
% Hebrew font — switched via ucharclasses (XeTeX) or
% directly referenced where needed
% \newfontfamily\hebrewfont{Alef}[
%   Script = Hebrew
% ]
% \newfontfamily\babelfont{Alef}[
%   Script = Hebrew
% ]
% \usepackage[hebrew, english]{babel}
% \babelfont[hebrew]{rm}{Alef}
% \usepackage[english, bidi=basic]{babel}
% \usepackage[english, hebrew, bidi=basic]{babel}
\usepackage[english, bidi=basic]{babel}

% \babelprovide[onchar=fonts ids, main=false]{hebrew} % all RTL, IPA wrongs
% \babelprovide[onchar=ids fonts, main=false]{english}
% \babelprovide[onchar=fonts ids, main=true]{english}
\babelprovide[onchar=fonts ids]{hebrew}

\babelfont[hebrew]{rm}{Alef}
% IPA font — called explicitly by ipa.lua filter output
\newfontfamily\ipafont{Source Sans 3}

% ------------------------------------------------------------
% Hebrew bidi — luatexja-style or direct LuaTeX primitive.
% Since babel is not loaded, we use the LuaTeX direction
% primitive directly: TRT = right-to-left, TLT = left-to-right
% Hebrew Unicode chars get RTL from the UBA via luaotfload.
% ------------------------------------------------------------
\directlua{
  luatexbase.add_to_callback("pre_linebreak_filter",
    function(head) return head end,
    "dummy_bidi")
}

\usepackage[hyperfootnotes=false]{hyperref}
% ============================================================
% header.tex  —  include-in-header snippet
% Replaces $var$ Pandoc syntax with \ifdefined\doc... guards.
% Vars injected via vars.lua Lua filter from YAML front matter:
%   \doctitle      \docsubtitle   \docauthor   \docdate
%   \docmeta       \docfoot       \doclogoleft \doclogoright
%
% No fancyhdr. Header fires via \AtBeginDocument,
% footer via \AtEndDocument — bidi-safe.
% ============================================================

\usepackage{graphicx}
\usepackage{tabularx}
\usepackage{xcolor}
\usepackage{array}
\usepackage{tikz}

% ---------- header font (adjust to taste) ----------
\newfontfamily\headerfont{Source Sans 3}

% ---------- footer colors ----------
\definecolor{footerbg}{HTML}{000000}
\definecolor{footerfg}{HTML}{FFFFFF}

% ---------- column types ----------
\newcolumntype{L}{>{\raggedright\arraybackslash}m{10mm}}
\newcolumntype{C}{>{\centering\arraybackslash}X}
\newcolumntype{R}{>{\raggedleft\arraybackslash}m{45mm}}

% vertical separator rule
\newcommand{\vsep}{\color{black}\vrule width 1pt}

% ============================================================
% CUSTOM HEADER
% ============================================================
\newcommand{\CustomHeader}{%
% \LTR{%
\noindent
\begin{tabularx}{\textwidth}{L !{\vsep} C !{\vsep} R}

% LEFT LOGO
\ifdefined\doclogoleft
  \includegraphics[width=10mm,height=10mm,keepaspectratio]{\doclogoleft}%
\fi
&
% CENTER: title + subtitle stacked in a minipage so \\ is a
% line break within the cell, not a new table row
\begin{minipage}[c]{\linewidth}
\centering
\ifdefined\doctitle
  {\Large\bfseries\headerfont\doctitle}%
\fi
\ifdefined\docsubtitle
  \\[2pt]{\normalsize\headerfont\docsubtitle}%
\fi
\end{minipage}
% &
% % CENTER TITLE BLOCK
% \ifdefined\doctitle
%   {\Large\bfseries\doctitle}%
% \fi
% \ifdefined\docsubtitle
%   \\\normalsize\docsubtitle
% \fi
&
% RIGHT LOGO
\ifdefined\doclogoright
  \includegraphics[height=10mm,keepaspectratio]{\doclogoright}%
\fi

\end{tabularx}

\hrule
% \vspace{6pt}

% ---------- META AREA ----------
{\small
\ifdefined\docauthor
  \textbf{Author:} \docauthor\\
\fi
\ifdefined\docdate
  \textbf{Date:} \docdate
\fi
\ifdefined\docmeta
  \\\docmeta\\
  snc: 16132.  
  {\ipafont ɪ ʃ ɔ ː ʁ}  
  {\ipafont wer es [pɛʟɛχɪʟ] ausspricht, und nicht [pɛʀɛχɪʟ]}
\fi
}
\vspace{6pt}

\hrule

% \vspace{12pt}
% }% end \LTR
}% end \CustomHeader

% ============================================================
% CUSTOM FOOTER  (overlay, anchored to page bottom)
% ============================================================
\newcommand{\CustomFooter}{%
\begin{tikzpicture}[remember picture, overlay]
\node[
  anchor=south,
  draw,
  fill=footerbg,
  text=footerfg,
  inner sep=10pt,
  rounded corners=2pt,
  text width=\dimexpr\textwidth-20pt\relax
]
at ([yshift=5mm]current page.south)
{%
  \begin{center}
  {\small
  \ifdefined\docfoot
    % {\docfoot}%
    {\headerfont\docfoot}%
  \fi
  }
  \end{center}
};
\end{tikzpicture}%
}% end \CustomFooter

% ============================================================
% FIRE header at doc start, footer at doc end
% ============================================================
% \AtBeginDocument{\CustomHeader}
\AtEndDocument{\CustomFooter}
\renewcommand{\maketitle}{}
\KOMAoption{captions}{tableheading}
\usepackage{fancyvrb}
\usepackage{fvextra}
\RecustomVerbatimEnvironment{verbatim}{Verbatim}{breaklines=true,breakanywhere=true,fontsize=\small}
\newcommand{\doctitle}{yakura replique}
\newcommand{\docdate}{2026-04-01}
\newcommand{\docmeta}{class: [16838] germanische sprachen im vergleich tutor: hüning/konvicka term: WS25/26}
\newcommand{\doclogoleft}{../assets/fadenlogo-fwb.png}
\newcommand{\doclogoright}{../assets/dcl-head-niemand.png}
\newcommand{\docauthor}{esteeschwarz}
\makeatletter
\@ifpackageloaded{caption}{}{\usepackage{caption}}
\AtBeginDocument{%
\ifdefined\contentsname
  \renewcommand*\contentsname{Table of contents}
\else
  \newcommand\contentsname{Table of contents}
\fi
\ifdefined\listfigurename
  \renewcommand*\listfigurename{List of Figures}
\else
  \newcommand\listfigurename{List of Figures}
\fi
\ifdefined\listtablename
  \renewcommand*\listtablename{List of Tables}
\else
  \newcommand\listtablename{List of Tables}
\fi
\ifdefined\figurename
  \renewcommand*\figurename{Figure}
\else
  \newcommand\figurename{Figure}
\fi
\ifdefined\tablename
  \renewcommand*\tablename{Table}
\else
  \newcommand\tablename{Table}
\fi
}
\@ifpackageloaded{float}{}{\usepackage{float}}
\floatstyle{ruled}
\@ifundefined{c@chapter}{\newfloat{codelisting}{h}{lop}}{\newfloat{codelisting}{h}{lop}[chapter]}
\floatname{codelisting}{Listing}
\newcommand*\listoflistings{\listof{codelisting}{List of Listings}}
\makeatother
\makeatletter
\makeatother
\makeatletter
\@ifpackageloaded{caption}{}{\usepackage{caption}}
\@ifpackageloaded{subcaption}{}{\usepackage{subcaption}}
\makeatother
\usepackage{bookmark}
\IfFileExists{xurl.sty}{\usepackage{xurl}}{} % add URL line breaks if available
\urlstyle{same}
\hypersetup{
  pdfauthor={esteeschwarz},
  colorlinks=true,
  linkcolor={blue},
  filecolor={Maroon},
  citecolor={Blue},
  urlcolor={Blue},
  pdfcreator={LaTeX via pandoc}}


\title{yakura replique}
\author{esteeschwarz}
\date{2026-04-01}
\begin{document}
\maketitle

\CustomHeader


\subsection{notes on replication
study}\label{notes-on-replication-study}

we started replicating the workflow presented in Yakura et al. (2025)
study beginning with creating the target corpus. Yakura et al. (2025)
built corpora from transcribed youtube video audios and podcasts,
referencing a research organisations registry (ROR (2025)) to collect
videos from institutional educational youtube channels. we assemlbled
data the same way restricted to german institutions. Yakura et al.
(2025) limited the target corpus to 4 years before the introduction of
the GPT chat agent on 2022-11-31 up to 2024-05-31. this allowed a
timebased analysis of hypothesised appearance of AI-speech induced
variances.

we tested the workflow until state of video transcribed using the
youtube search API for finding the corresponding video channels,
llama3.2 to match the correct channel within the search results, yt-dlp
to download the video audio, ffmpeg library to convert to PCM .wav and
whisper AI to finally transcribe the audio. all worked well
(\href{https://github.com/esteeschwarz/SPUND-LX/tree/main/germanic/HA/LLM-001.R}{script})
with resulting two texts from Mannheim University channel youtube
contributions. we estimate an overall server runtime of about 10h to
download and transcribe to text all audio of above categorized channels.

\subsection{process on yakura
replication}\label{process-on-yakura-replication}

after starting with the same pipeline of building a corpus from youtube
material, we decided for another alternative out of resource reasons.
the current corpus which we will work on is created from german
parliamental protocols, freely available (resource reason) here: DIP
(2026). Also we decided for google gemini since working with that model
simply costs less than the openAI GPT variant.

for the beginning we created a subset of protocols from 2021-01-01 to
2021-07-31 which are 38 protocols. to get the model preferred vocabulary
we prompted gemini to summarize each protocol in its own words with
restricting to not using more than 5\% of words from the original text
and limited to 300 words/summary. See prompt text
Section~\ref{sec-gemini-prompt}. we postag the corpus and devise
relative lemma frequencies to get the gemini keywords.

\subsubsection{gemini prompt}\label{sec-gemini-prompt}

\begin{verbatim}
[1] "System prompt: "                                                                                                                                                                                                                                                                                                                                                                                                                                                                                                                                                                           
[2] "You are a member of german parliament. Prepare a summary of the text provided to present at a local community meeting of your party members. Output in german language, no preamble, no extra information, just the plain text. Wordcount maximal 300 words, containing not more than 5% of the keywords of the text provided and explicitly not just a list of keywords but an entertaining text. You are supposed to interprete freely, including background insights on daily politics. Keep in mind thatthe text will be used as is as keynotes to the talk being held to the locals. "
[3] "Text:"                                                                                                                                                                                                                                                                                                                                                                                                                                                                                                                                                                                     
\end{verbatim}

\begin{quote}
please preview/follow paper \href{gemini.html}{here}.
\end{quote}

\phantomsection\label{refs}
\begin{CSLReferences}{1}{0}
\bibitem[\citeproctext]{ref-dip_dip_2026}
DIP. 2026. {``{DIP} - {Bundestagsprotokolle}.''} Docs. \emph{DIP - API}.
Berlin.
\url{https://dip.bundestag.de/\%C3\%BCber-dip/hilfe/api\#content}.

\bibitem[\citeproctext]{ref-ror_ror_2025}
ROR, Research Organization Registry. 2025. {``{ROR} {Data}.''} Zenodo.
\url{https://doi.org/10.5281/zenodo.17953395}.

\bibitem[\citeproctext]{ref-yakura_empirical_2025}
Yakura, Hiromu, Ezequiel Lopez-Lopez, Levin Brinkmann, Ignacio Serna,
Prateek Gupta, Ivan Soraperra, and Iyad Rahwan. 2025. {``Empirical
Evidence of {Large} {Language} {Model}'s Influence on Human Spoken
Communication.''} arXiv.
\url{https://doi.org/10.48550/arXiv.2409.01754}.

\end{CSLReferences}




\end{document}
